\subsubsection{女性性少数者的社群与支持资源：一份可用的清单}

前几节介绍了全球性健康机构与通用性健康资源（见《性健康资源与支持》一章）。本节补上一份更贴近实际使用的清单——按"需要什么"来组织，而不是按机构类型。

\textbf{一、检测与预防资源}
\begin{itemize}
\item \textbf{疾控中心（CDC）自愿咨询检测门诊（VCT）}：多数城市疾控中心及其指定的社区组织提供 HIV、梅毒等项目的免费或低成本检测，部分可匿名；
\item \textbf{自检试剂}：HIV 抗体自检试剂可通过正规渠道购买（认准药品或医疗器械批准文号），它适合作为初筛，\textbf{阳性结果必须到疾控或医院复核确认}；
\item \textbf{PrEP（暴露前预防）与 PEP（暴露后阻断）}：国内部分城市的定点医院与疾控已开展 PrEP 服务，可咨询当地疾控是否属于可及区域；PEP 需在 72 小时内启动，属于急诊级事项；
\item \textbf{HPV 疫苗与宫颈筛查}：社区与妇幼保健院均可预约；有性生活的女性按指南定期筛查，不因性伴侣性别而降低频次。
\end{itemize}

\textbf{二、友善医疗与心理支持}
\begin{itemize}
\item \textbf{LGBTQ+ 友善医疗机构与医生}：部分城市的综合医院、妇幼保健院或民营性健康门诊有友善项目。判断标准可以很实际——问一句"你们接待女性同性伴侣就诊吗"，对方的反应本身就是答案；
\item \textbf{心理咨询}：需要区分两类资源，一是面向性少数群体的友善心理咨询师（可通过社群组织获取推荐名单）；二是危机干预热线，用于急性心理危机（情绪崩溃、自伤念头等），全国与各省市均有 24 小时心理援助热线，公开可查；
\item \textbf{性侵受害者援助}：遭遇性侵后可到指定的性侵取证医疗机构进行取证与医学处置（72 小时内评估 PEP 与紧急避孕），同时可联系妇女维权热线与法律援助机构；是否报警由当事人决定，但取证与医学处理不建议等待。
\end{itemize}

\textbf{三、社群组织与同伴支持}
\begin{itemize}
\item \textbf{社群组织的价值}：由社群自己运营的组织通常能提供医院清单、体检指引、法律咨询转介与同伴经验分享——信息不一定都能在公开渠道找到，因为它们的价值恰恰在于"过来人知道怎么走"；
\item \textbf{找到它们的可靠方式}：通过本地 LGBTQ+ 社群、高校相关社团、社群公益组织的公开账号；参加活动时对信息边界保持谨慎（\textbf{参加活动不等于同意公开身份}，社群内部同样有隐私保护规则）；
\item \textbf{线上社区}：论坛、社群与内容平台上的自助小组能提供日常陪伴与经验问答。需注意：线上一是信息真伪混杂，二是存在借社群名义的诈骗（收费"内部资源""包过体检"多不可信）。判断标准很简单——正规公益服务不向求助者收取高额费用。
\end{itemize}

\textbf{四、学习与自助}
\begin{itemize}
\item \textbf{书籍}：除本卷前面推荐的性学著作外，可以留意带有"女性性体验"视角的作品——本书《性与法律：同意、权益与边界》章与《LGBTQ+人群的性健康与权益》章中也整理了相关经验叙述；
\item \textbf{给伴侣的功课}：如果伴侣需要了解你的处境，与其转述二手信息，不如一起读一节、看一份公开指南——\textbf{共同学习比单方面解释更容易建立理解}，这也符合本书反复强调的沟通原则；
\item \textbf{自助不等于独自硬扛}：任何持续两周以上的低落、焦虑、睡眠受影响或出现身体症状的情况，都值得走一遍专业渠道。本章前几节已给出就医路径，这里不再重复。
\end{itemize}

\noindent\textit{【编者按】}本节不提供具体机构的名称与联系方式：一来各地资源差异大且时常变动，二来这类信息需要以最新官方发布为准。可执行的做法是记住上面五类渠道，然后通过当地疾控、社群公开账号或医院官方渠道核实到具体机构。
